\documentclass{beamer}
\usepackage{beamer-custom}

\usefonttheme{serif}
% Theme
\usetheme{Montpellier}
\usecolortheme{beaver}
\useinnertheme{circles}

% Navigation
\setbeamertemplate{navigation symbols}{}
\setbeamertemplate{footline}[frame number]

% Frame title formatting
\setbeamertemplate{frametitle}{
    \vspace{0.5em}
    \insertframetitle
    \par
}
%\usepackage{cmbright,bbold}
%\renewcommand{\mathcal}{\mathtt}
%\usepackage{kpfonts}

\AtBeginSection[]
{
  \begin{frame}
    \frametitle{Contents}
    \tableofcontents[currentsection]
  \end{frame}
}
\usepackage[backend=biber, style=alphabetic,sorting=nyt]{biblatex}
\addbibresource{references.bib}

\title{The Free Central Limit Theorem}
\author{Avinash Iyer}
\institute{University of Virginia}
\date{April 6, 2026}
\begin{document}
\begin{frame}
  \titlepage
\end{frame}

\begin{frame}
  \frametitle{Outline}
  \tableofcontents
\end{frame}
\section{Definitions}%
\begin{frame}[allowframebreaks]
  \frametitle{Noncommutative Probability}
  \begin{itemize}
    \item A \textbf{$C^{\ast}$-probability space} is a $C^{\ast}$-algebra is a pair $\left( \mathcal{A},\varphi \right)$, where $\mathcal{A}$ is a (unital) $C^{\ast}$-algebra --- a complete normed $\ast$-algebra satisfying $\norm{a^{\ast}a} = \norm{a}^2$ --- and an injective linear functional $\varphi\colon \mathcal{A}\rightarrow \C$ satisfying $\varphi(1) = 1$ and $\varphi\left( a^{\ast}a \right)\geq 0$ for all $a\in A$, with equality only if $a = 0$ (such a $\varphi$ is called a state). We call elements $a\in \mathcal{A}$ \textit{noncommutative random variables}.
    \item If $\set{\mathcal{A}_i}_{i\in I}\subseteq \mathcal{A}$ is a family of $C^{\ast}$-subalgebras, the family is said to be \textit{freely independent} if $\varphi\left( a_1\cdots a_k \right) = 0$ whenever
      \begin{itemize}
        \item $k > 0$;
        \item $a_j\in \mathcal{A}_{i_j}$ for each $j$;
        \item $\varphi\left( a_j \right) = 0$ for each $j$;
        \item $i_1\neq i_2$, $i_2\neq i_3$, etc.
      \end{itemize}
      If the $\mathcal{A}_i = C^{\ast}\left( a_i \right)$ are generated by single elements, then the elements $\set{a_i}_{i\in I}$ are called freely independent.
    \item If $a\in \mathcal{A}$ is an element of a $C^{\ast}$-probability space, the \textit{moments} of $a$ are the values $\set{\varphi\left( a^n \right) | n=0,1,\dots}$.
    \item If $\left( \mathcal{A}_N,\varphi_N \right)$ is a sequence of $C^{\ast}$-probability spaces, and $\left( \mathcal{A},\varphi \right)$ is another $C^{\ast}$-probability space, then a sequence $\left( a_N \right)_N$ with each $a_N\in \mathcal{A}_N$ is said to \textit{converge in distribution} to $a\in \mathcal{A}$ if, for each $k\in \N$,
      \begin{align*}
        \lim_{N\rightarrow\infty} \varphi_N\left( a_N^{k} \right) &= \varphi\left( a^k \right).
      \end{align*}
  \end{itemize}
\end{frame}
\begin{frame}[allowframebreaks]
  \frametitle{Combinatorics}

\end{frame}
\begin{frame}[allowframebreaks]
  \frametitle{References}
  \nocite{*}
  {\tiny \printbibliography}
\end{frame}
\end{document}
